%------------------------------------------
% Considerações Finais
%------------------------------------------
\section{Considerações Finais}
\label{sec:conclusao}


\todo[inline]{\textbf{NOTA}: No texto de qualificação, as considerações finais têm caráter parcial e devem permitir que a banca avalie o andamento da pesquisa e a viabilidade de concluí-la no prazo regimental do mestrado. Descreva, de forma resumida: (1) o que já foi desenvolvido, relacionando-o aos objetivos específicos; (2) os resultados preliminares e as contribuições esperadas, tanto científicas quanto o produto técnico ou tecnológico e o seu impacto no contexto de aplicação; (3) os riscos e os desafios para as próximas etapas, com as respectivas estratégias de mitigação; e (4) o cronograma das atividades até a defesa da dissertação. Exemplo abaixo:}

\textcolor{red}{O objetivo deste estudo é [PREENCHER]. Para alcançá-lo, adota-se [PREENCHER COM UMA SÍNTESE DE DUAS A TRÊS LINHAS SOBRE O MÉTODO]. Até o momento, foram concluídas as seguintes etapas: [PREENCHER COM AS ATIVIDADES REALIZADAS E OS OBJETIVOS ESPECÍFICOS JÁ ATENDIDOS]. Os resultados preliminares indicam que [PREENCHER COM A SÍNTESE DOS RESULTADOS PRELIMINARES E COM O QUE ELES SUGEREM EM RELAÇÃO À PERGUNTA DE PESQUISA].}

\textcolor{red}{Como contribuições esperadas, destacam-se [PREENCHER COM AS CONTRIBUIÇÕES CIENTÍFICAS] e o desenvolvimento de [PREENCHER COM O PRODUTO TÉCNICO OU TECNOLÓGICO], que deverá beneficiar [PREENCHER COM A ORGANIZAÇÃO, O SETOR OU O PÚBLICO-ALVO]. Como próximas etapas, pretende-se [PREENCHER COM AS ATIVIDADES PENDENTES]. Os principais desafios identificados são [PREENCHER COM OS RISCOS, POR EXEMPLO, ACESSO A DADOS OU A USUÁRIOS DA ORGANIZAÇÃO PARCEIRA], que serão mitigados por meio de [PREENCHER COM AS ESTRATÉGIAS DE MITIGAÇÃO].}

\textcolor{red}{Pretende-se, ainda, submeter os resultados da pesquisa a [PREENCHER COM OS PERIÓDICOS OU CONFERÊNCIAS-ALVO] e registrar ou disponibilizar o produto desenvolvido por meio de [PREENCHER, POR EXEMPLO, REGISTRO DE SOFTWARE NO INPI OU REPOSITÓRIO PÚBLICO].}


\subsection{Cronograma}
\label{sub:cronograma}

\todo[inline]{\textbf{NOTA}: Apresente o cronograma das atividades a serem realizadas entre a qualificação e a defesa, compatível com o prazo regimental do PPCA. Ajuste as atividades e os meses ao seu trabalho. Segue um \textbf{exemplo}.}

A Tabela~\ref{tab:cronograma} apresenta o cronograma das atividades previstas até a defesa da dissertação.

\begin{table}[htbp]
  \centering
  \scalefont{0.9}
  \begin{tabular}{|l|c|c|c|c|c|c|}
    \hline
    \multicolumn{1}{|c|}{\multirow{2}{*}{\textbf{Atividades}}} & \multicolumn{6}{c|}{\textbf{Meses}} \bigstrut\\
    \cline{2-7}
    & \textbf{1} & \textbf{2} & \textbf{3} & \textbf{4} & \textbf{5} & \textbf{6} \bigstrut\\
    \hline
    Ajustes sugeridos pela banca de qualificação & X &  &  &  &  &  \bigstrut\\
    \hline
    Implementação do produto técnico ou tecnológico & X & X & X &  &  &  \bigstrut\\
    \hline
    Execução da avaliação experimental &  &  & X & X &  &  \bigstrut\\
    \hline
    Análise dos resultados &  &  &  & X & X &  \bigstrut\\
    \hline
    Submissão de artigo científico &  &  &  &  & X &  \bigstrut\\
    \hline
    Escrita da dissertação &  & X & X & X & X & X \bigstrut\\
    \hline
    Defesa da dissertação &  &  &  &  &  & X \bigstrut\\
    \hline
  \end{tabular}
  \caption{Cronograma das atividades até a defesa.}
  \label{tab:cronograma}
\end{table}
